% Figure: ablations, six Fund R hard runs, 193 errors, frozen v1.
% Source: out/hard_runs/ablate_frozen_summary.json (tier1_fixed_correctly / sent_to_tier2 /
% not_flagged / tier1_wrong_changes). no_pagecheck and no_notprinted equal "full" (189/2/2/0).
% Closer-not-exact errors are also in a tier 2 item, so fixed + tier 2 + missed = 193 in every row.
\begin{figure}[t]
\centering
\begin{tikzpicture}
\begin{axis}[qbar, width=3.7cm, xmin=0, xmax=193, y=13pt,
  xlabel={tier 0 errors (6 hard runs, $n=193$)},
  ytick={1,2,3,4,5,6}, ymin=1, ymax=6, y dir=reverse,
  yticklabels={Full system$^{*}$, No misread repair, No swap/period/sign search, Page only, No label read, Arithmetic only},
  xtick={0,50,100,150,193},
  legend columns=3, legend style={at={(0.5,1.03)}, anchor=south},
  clip=false]
\addplot[fixbar] coordinates {(189,1) (187,2) (181,3) (179,4) (14,5) (5,6)};
\addplot[t2bar] coordinates {(2,1) (4,2) (10,3) (12,4) (166,5) (124,6)};
\addplot[missbar] coordinates {(2,1) (2,2) (2,3) (2,4) (13,5) (64,6)};
\legend{fixed, to tier 2, missed}
% direct value labels for fixed
\coordinate (abl1) at (axis cs:94,1);
\coordinate (abl2) at (axis cs:93,2);
\coordinate (abl3) at (axis cs:90,3);
\coordinate (abl4) at (axis cs:89,4);
\coordinate (abl5) at (axis cs:97,5);
\coordinate (abl6) at (axis cs:67,6);
\coordinate (abl7) at (axis cs:161,6);
% wrong changes
\coordinate (abl8) at (axis cs:193,3);
\coordinate (abl9) at (axis cs:193,4);
\end{axis}
\node[inlab] at (abl1) {189};
\node[inlab] at (abl2) {187};
\node[inlab] at (abl3) {181};
\node[inlab] at (abl4) {179};
\node[darklab] at (abl5) {166};
\node[darklab] at (abl6) {124};
\node[misslab] at (abl7) {64};
\node[wronglab] at (abl8) {\textcolor{cWrong}{$\blacktriangle$}\,2 wrong};
\node[wronglab] at (abl9) {\textcolor{cWrong}{$\blacktriangle$}\,2 wrong};
\end{tikzpicture}
\caption{Ablations of PE v1's tier~1 on the six Fund~R hard runs. Without the search, 2 label reads with the wrong sign are applied ($\blacktriangle$). $^{*}$Turning off the cited-spot check or the ``not printed anywhere'' flag gives the same counts. Wilson intervals are in Table~\ref{tab:ablation}.}
\label{fig:ablation}
\end{figure}
